% !PRINTED 746
\weberpage{746}
%<<CONT
\gesperrt{oder einem Theil derselben in eine gegenseitig eindeutige Beziehung gesetzt werden kann}\footnote{\gesperrt{Cantor}, l.\,c. Der Begriff der abzählbaren Mengen deckt sich mit dem von \gesperrt{Dedekind} ohne die \gesperrt{Voraussetzung des Zahlensystems} definirten Begriffe der \gesperrt{einfach unendlichen Systeme}. (\gesperrt{Dedekind}, Was sind und was sollen die Zahlen? \S.\,6. Braunschweig 1887. Zweite unveränderte Auflage 1893.)}.

Jede endliche Menge ist hiernach abzählbar, und bei der Abzählung wird die natürliche Zahlenreihe nur bis zu einer gewissen höchsten Zahl verwendet. Im Weiteren betrachten wir vorzugsweise unendliche Mengen.

Wir können der Definition für eine unendliche abzählbare Menge auch den Ausdruck geben, dass es eine solche Menge ist, bei der jedem Elemente eine bestimmte Zahl der natürlichen Zahlenreihe als Name beigelegt werden kann, so dass auch jede Zahl der Zahlenreihe dabei verwandt wird, also eine Menge, in der es ein erstes, zweites, drittes \ldots hundertstes \ldots Element giebt.

Endlich können wir auch sagen, dass eine unendliche abzählbare Menge eine solche ist, die sich derart in eine Reihe ordnen lässt, dass ein \gesperrt{erstes} Element vorhanden ist, und dass auf jedes Element ein \gesperrt{bestimmtes anderes der Menge} folgt, und jedem Elemente, mit Ausnahme des ersten, ein \gesperrt{bestimmtes anderes Element vorangeht}. Eine solche Reihe können wir eine zählbare Anordnung nennen.

Es ist klar, dass eine abzählbare Menge nicht nur auf eine, sondern auf unendlich viele verschiedene Arten abzählbar ist.

Ausser der natürlichen Zahlenreihe selbst, die sicher abzählbar ist, können wir als zweites Beispiel das System der rationalen positiven echten Brüche anführen, die unter Anderem in folgender Weise abgezählt werden können:
\[
\frac{1}{2}, \frac{1}{3}, \frac{2}{3}, \frac{1}{4}, \frac{3}{4}, \frac{1}{5}, \frac{2}{5}, \frac{3}{5}, \frac{4}{5}, \frac{1}{6}, \frac{5}{6}, \ldots
\]
d.\,h. so, dass jeder grössere Nenner dem kleineren Nenner folgt, und dass die Brüche von gleichem Nenner nach der Grösse des Zählers geordnet sind. Wollte man aber die Brüche nach der Grösse ihres numerischen Werthes ordnen, so würde man keine zählbare Anordnung bekommen.

2. \gesperrt{Die Gesammtheit aller algebraischen Zahlen ist eine abzählbare Menge.}
